% LaTeX source maintained directly for Overleaf.
\chapter*{Lời mở đầu}
\addcontentsline{toc}{chapter}{Lời mở đầu}

Knowledge Graph là một hướng tiếp cận phù hợp với dữ liệu có nhiều loại thực thể
và quan hệ liên kết. Trong miền phim, một bộ phim liên hệ với diễn viên, đạo diễn,
thể loại, từ khóa và hãng sản xuất; các quan hệ này tạo ra nhiều câu hỏi mà cách
biểu diễn đồ thị có thể mô tả trực tiếp và dễ theo dõi hơn.

Báo cáo trình bày quá trình xây dựng Movie Knowledge Graph từ dữ liệu TMDB và
IMDb, bao gồm thu thập và chuẩn hóa dữ liệu, thiết kế ontology và graph schema,
nạp dữ liệu vào Neo4j, truy vấn Cypher, biểu diễn RDF/OWL và một số luật suy diễn.
Trên nền tảng đó, hệ thống triển khai hai chức năng chính: hỏi–đáp về phim bằng
ngôn ngữ tự nhiên và gợi ý phim kèm lý do dựa trên các quan hệ chung trong graph.

Mục tiêu của đề tài là xây dựng và đánh giá một quy trình Knowledge Graph đầu
cuối có thể vận hành và kiểm chứng. Trọng tâm nghiên cứu gồm tích hợp dữ liệu đa
nguồn với định danh ổn định, truy vấn nhiều bước, suy diễn có bằng chứng, hỏi–đáp
an toàn và gợi ý phim có khả năng giải thích. Các kết quả được gắn với snapshot,
cấu hình và protocol cụ thể để phân biệt bằng chứng thực nghiệm với nhận định
khái quát.

Nội dung báo cáo được tổ chức thành bảy chương:

\begin{enumerate}
\item \textbf{Chương 1 – Đặt vấn đề và tổng quan:} trình bày bối cảnh, bài toán,
câu hỏi nghiên cứu, phạm vi, các công trình liên quan và lý do lựa chọn mô hình
Knowledge Graph.
\item \textbf{Chương 2 – Lý thuyết và mô hình dữ liệu:} giới thiệu nền tảng về
biểu diễn tri thức, RDF/RDFS/OWL, Property Graph, cơ chế suy diễn, entity
resolution và provenance.
\item \textbf{Chương 3 – Thiết kế và cài đặt hệ thống:} mô tả kiến trúc, nguồn
dữ liệu, pipeline xử lý, ontology, lược đồ Neo4j, RDF export và cấu hình các
engine suy diễn.
\item \textbf{Chương 4 – Truy vấn, xử lý và nghiệp vụ:} trình bày các truy vấn
Cypher/SPARQL, competency question, cơ chế bảo đảm an toàn truy vấn và các luật
suy diễn có bằng chứng.
\item \textbf{Chương 5 – Thực nghiệm và đánh giá:} mô tả dữ liệu, protocol,
baseline, chỉ số đánh giá và kết quả về chất lượng graph, entity resolution,
reasoning, hỏi–đáp, recommendation và hiệu năng.
\item \textbf{Chương 6 – Ứng dụng và vận hành hệ thống:} trình bày các chức năng
hỏi–đáp, gợi ý phim, API, giao diện, cơ chế giải thích và các tình huống lỗi.
\item \textbf{Chương 7 – Kết luận và hướng phát triển:} tổng hợp kết quả, trả lời
câu hỏi nghiên cứu, nêu hạn chế và đề xuất các hướng mở rộng.
\end{enumerate}

Phần phụ lục cung cấp hướng dẫn cài đặt và chạy chương trình, ví dụ sử dụng API,
một số truy vấn minh họa và vị trí các tệp dữ liệu, kết quả thực nghiệm.

\bigskip\hrule\bigskip
